\documentclass[11pt]{article}

\usepackage[a4paper,margin=28mm]{geometry}
\usepackage{amsmath,amssymb,amsthm,mathtools}
\usepackage{booktabs}
\usepackage{enumitem}
\usepackage{microtype}
\usepackage[hidelinks]{hyperref}
\usepackage{xcolor}

\makeatletter
\g@addto@macro\UrlBreaks{\do\.\do\_\do\;\do\-}
\makeatother
\hypersetup{
  pdftitle={Verified Statements for a Finite-Precision Lanczos Analysis},
  pdfauthor={Anonymous formalization report},
  pdfsubject={Journal-style statement record for a Lean 4 formalization}
}

\definecolor{leanblue}{RGB}{24,72,120}
\newcommand{\R}{\mathbb{R}}
\newcommand{\N}{\mathbb{N}}
\newcommand{\norm}[1]{\lVert #1\rVert_2}
\newcommand{\abs}[1]{\lvert #1\rvert}
\newcommand{\dist}{\operatorname{dist}}
\newcommand{\spec}{\operatorname{spec}}
\newcommand{\diag}{\operatorname{diag}}
\newcommand{\supp}{\operatorname{supp}}
\newcommand{\Wone}{\mathcal W_1^{\mathrm{dual}}}
\newcommand{\lean}[1]{\texorpdfstring{\textcolor{leanblue}{\nolinkurl{#1}}}{#1}}
\newcommand{\formalnote}[1]{%
  \par\smallskip\noindent{\small\emph{Formal declaration(s):}}%
  \par\noindent{\scriptsize #1.}%
  \par\smallskip}
\setlength{\emergencystretch}{3em}

\newtheorem{assumption}{Assumption}
\newtheorem{theorem}{Theorem}[section]
\newtheorem{corollary}[theorem]{Corollary}
\newtheorem{proposition}[theorem]{Proposition}
\theoremstyle{definition}
\newtheorem{definition}[theorem]{Definition}
\newtheorem{remark}[theorem]{Remark}

\title{Verified Statements for a Finite-Precision Lanczos Analysis\\
\large A Formal Statement Record}
\author{Anonymous formalization report}
\date{20 August 2026}

\begin{document}

\maketitle

\begin{abstract}
This document gives a mathematical rendering of the principal statements
checked in a Lean~4 development of finite-precision Lanczos analysis.  The
formal input is not an implementation of floating-point arithmetic.  It is a
real symmetric matrix together with computed vectors and coefficients that
satisfy explicitly stated local arithmetic identities and error bounds.  From
that interface the development verifies explicit local-error estimates, a
Paige-type commutator identity and loss-of-orthogonality bound, Ritz-value
containment and stabilization, and a Greenbaum-type backward model.  The
backward model is a symmetric tridiagonal matrix whose leading block is the
computed Lanczos tridiagonal, which reproduces that block under exact Lanczos,
and whose spectrum and starting spectral measure are controlled relative to
the original matrix.  Every constant and hypothesis used below is displayed.
The purpose is auditability: the theorem names attached to the statements
identify the corresponding Lean declarations.
\end{abstract}

\section{Purpose, scope, and formal status}

The authoritative objects are the Lean sources in the accompanying project.
This manuscript translates their paper-facing declarations into conventional
numerical-linear-algebra notation.  It is deliberately conservative: no
asymptotic notation is substituted for a proved constant, and conclusions
that are conditional in Lean remain conditional here.

The project builds with Lean~4 and Mathlib using the pinned toolchain
\texttt{v4.32.2}.  A full \texttt{lake build} completed successfully.  An
axiom audit of the principal results reported only \texttt{propext},
\texttt{Classical.choice}, and \texttt{Quot.sound}, the standard logical
axioms used by Lean and Mathlib.  A source search found no occurrences of
\texttt{sorry}, \texttt{admit}, or user-declared axioms.

\begin{center}
\fbox{\begin{minipage}{0.91\textwidth}
\textbf{Central modelling limitation.}
The Lean development does not prove that executing a particular machine-level
Algorithm~1 produces the formal run described below.  It proves that
\emph{any} family of real objects satisfying Assumption~\ref{ass:run} obeys
the conclusions in this document.  Connecting an implementation, a
floating-point semantics, or a probabilistic rounding model to that interface
is a separate obligation.
\end{minipage}}
\end{center}

\section{Notation}

All matrices and vectors are real.  The vector norm and matrix norm are the
Euclidean norm and its induced operator $2$-norm, respectively.  Transpose is
written ${}^{T}$, and $e_j$ denotes a coordinate vector.  The indices in the
mathematical presentation are one-based.  The formal development uses natural
indices with the phantom values $v_0=0$ and $\beta_0=0$.

For a real symmetric matrix $H\in\R^{n\times n}$, write
\[
  \lambda_{\min}(H)=\min\spec(H),\qquad
  \lambda_{\max}(H)=\max\spec(H),
  \qquad
  d_H(\theta)=\dist(\theta,\spec(H)).
\]
The formal definitions return zero outside the symmetric, positive-dimensional
case; Assumption~\ref{ass:run} implies symmetry, and the run assumptions imply
$n>0$, so only the usual spectral interpretation occurs in the results below.

For $x\in\R^n$ and a symmetric $A$, define its starting spectral mass by
\begin{equation}\label{eq:spectral-measure}
  \mu_{A,x}=\sum_i (u_i^Tx)^2\,\delta_{\lambda_i(A)},
\end{equation}
where $(u_i)$ is the orthonormal eigenbasis selected by the formal spectral
theorem.  Repeated eigenvalues are combined because the formal measure is a
finitely supported function on $\R$.

For finite atomic masses $\mu$ and $\nu$, the development defines
\begin{equation}\label{eq:dual-wasserstein}
 \Wone(\mu,\nu)
 =\sup_{\operatorname{Lip}(f)\leq 1}
   \abs{\int f\,d\mu-\int f\,d\nu}.
\end{equation}
This is the dual quantity named \lean{atomicWasserstein} in the code.  We retain
the superscript ``dual'' here to make the formal definition explicit.

\section{The formal finite-precision interface}

\begin{assumption}[Lanczos run interface]\label{ass:run}
Let $n,k\in\N$.  The input to every run-level result consists of
\[
 H\in\R^{n\times n},\quad
 v_j,u_j,w_j,z_j,e^u_j,e^w_j,e^z_j,\Delta_j\in\R^n,
 \quad
 \alpha_j,\beta_j,e^\alpha_j\in\R,
 \quad \varepsilon\in\R,
\]
subject to all of the following conditions.

\begin{enumerate}[label=(A\arabic*),leftmargin=3.2em]
\item $H^T=H$, $\varepsilon\geq0$, $k>0$, and
      $k\varepsilon\leq 1/100$.
\item The phantom entries satisfy $v_0=0$ and $\beta_0=0$.
\item For every $j=0,\ldots,k-1$,
\begin{align}
 u_{j+1}&=Hv_{j+1}+e^u_{j+1},
 &\norm{e^u_{j+1}}
   &\leq\varepsilon\norm H\norm{v_{j+1}},\label{eq:run-u}\\
 w_{j+1}&=u_{j+1}-\beta_jv_j+e^w_{j+1},
 &\norm{e^w_{j+1}}
   &\leq\varepsilon\bigl(\norm{u_{j+1}}
       +\abs{\beta_j}\norm{v_j}\bigr),\label{eq:run-w}\\
 \alpha_{j+1}&=v_{j+1}^Tw_{j+1}+e^\alpha_{j+1},
 &\abs{e^\alpha_{j+1}}
   &\leq\varepsilon\norm{v_{j+1}}\norm{w_{j+1}},\label{eq:run-a}\\
 z_{j+1}&=w_{j+1}-\alpha_{j+1}v_{j+1}+e^z_{j+1},
 &\norm{e^z_{j+1}}
   &\leq\varepsilon\bigl(\norm{w_{j+1}}
      +\abs{\alpha_{j+1}}\norm{v_{j+1}}\bigr),\label{eq:run-z}\\
 \beta_{j+1}v_{j+2}&=z_{j+1}+\Delta_{j+1},
 &\norm{\Delta_{j+1}}&\leq\varepsilon\norm{z_{j+1}}.
 \label{eq:run-b}
\end{align}
\item The computed Lanczos vectors are nearly normalized:
\[
 \abs{\norm{v_j}^2-1}\leq\varepsilon,
 \qquad 1\leq j\leq k.
\]
If $\beta_k\neq0$, the additional vector satisfies
$\abs{\norm{v_{k+1}}^2-1}\leq\varepsilon$.
\item $\beta_j\geq0$ for $0\leq j\leq k$, and
$\beta_j>0$ for $1\leq j<k$.
\end{enumerate}
\end{assumption}
\formalnote{\lean{FinitePrecisionLanczos.LanczosRun}}

There are no other numerical assumptions at the run interface.  In
particular, $k\leq n$ is not assumed, $\beta_k$ need not be positive, and no
orthogonality between distinct $v_i$ and $v_j$ is assumed.  The condition
$\beta_j>0$ only excludes an internal exact breakdown before the last
reported step.

Define, for $1\leq j\leq k$,
\begin{align}
 f_j&=Hv_j-\beta_{j-1}v_{j-1}-\alpha_jv_j-\beta_jv_{j+1},\label{eq:def-f}\\
 g_j&=v_j^Tv_j-1,\qquad
 p_j=v_j^T(\beta_jv_{j+1}),
 \quad p_0=0.\label{eq:def-gp}
\end{align}
Let
\[
 V=[v_1\ \cdots\ v_k],\qquad F=[f_1\ \cdots\ f_k],
\]
and let $T_k$ be the symmetric tridiagonal matrix with diagonal
$(\alpha_1,\ldots,\alpha_k)$ and sub/superdiagonal
$(\beta_1,\ldots,\beta_{k-1})$.  Finally set
\[
 b=\beta_kv_{k+1},\qquad B=b e_k^T.
\]

\section{Local errors and assembled recurrence}

\begin{theorem}[Explicit local errors]\label{thm:local}
Under Assumption~\ref{ass:run}, for every $1\leq j\leq k$,
\begin{equation}\label{eq:local-bounds}
 \norm{f_j}\leq 700j\,\varepsilon\norm H,
 \qquad
 \abs{g_j}\leq\varepsilon,
 \qquad
 \abs{p_j}\leq1200j\,\varepsilon\norm H.
\end{equation}
Moreover,
\begin{equation}\label{eq:assembled}
  HV=VT_k+B+F,
  \qquad
  \norm F\leq700k^2\varepsilon\norm H.
\end{equation}
\end{theorem}
\formalnote{\lean{LanczosRun.local_errors}; \lean{LanczosRun.governing_matrix}; \lean{LanczosRun.Fmat_bound}}

For audit purposes, the local proof also verifies the componentwise estimates
\begin{equation}\label{eq:component-errors}
 \norm{e^u_j}\leq2j\varepsilon\norm H,
 \quad \norm{e^w_j}\leq200j\varepsilon\norm H,
 \quad \abs{e^\alpha_j}\leq200j\varepsilon\norm H,
\end{equation}
and
\begin{equation}\label{eq:component-errors-2}
 \norm{e^z_j}\leq300j\varepsilon\norm H,
 \qquad \norm{\Delta_j}\leq100j\varepsilon\norm H.
\end{equation}
These are the declarations \lean{LanczosRun.eu_explicit},
\lean{LanczosRun.ew_explicit}, \lean{LanczosRun.ealpha_explicit},
\lean{LanczosRun.ez_explicit}, and \lean{LanczosRun.delta_explicit}.

\section{Gram commutator and Paige bounds}

Set
\[
 \Omega=V^TV-I,
 \qquad D=\diag(g_1,\ldots,g_k),
\]
and let $R$ be the strictly upper-triangular part of $\Omega$.  Thus
\begin{equation}\label{eq:gram-split}
  \Omega=R^T+D+R.
\end{equation}
Introduce
\begin{align*}
 C_b&=V^TB,&
 G&=V^TF-F^TV,\\
 Q&=T_kD-DT_k,&
 N&=\diag(p_0-p_1,p_1-p_2,\ldots,p_{k-1}-p_k).
\end{align*}
For a matrix $M$, write $M_{\mathrm{up}}$ for its strictly upper-triangular
part.

\begin{theorem}[Gram commutator]\label{thm:commutator}
Under Assumption~\ref{ass:run},
\begin{equation}\label{eq:full-commutator}
 T_k\Omega-\Omega T_k=C_b-C_b^T+G,
\end{equation}
and every entry of $G$ satisfies
\begin{equation}\label{eq:G-bound}
 \abs{G_{rs}}\leq2800k\,\varepsilon\norm H.
\end{equation}
The upper-triangular commutator has the exact form
\begin{equation}\label{eq:upper-commutator}
 T_kR-RT_k=C_b+N+G_{\mathrm{up}}-Q_{\mathrm{up}}.
\end{equation}
Furthermore, $Q_{\mathrm{up}}$ is supported on the first superdiagonal and
\begin{equation}\label{eq:Q-bound}
 \abs{(Q_{\mathrm{up}})_{rs}}
 \leq600k\,\varepsilon\norm H.
\end{equation}
\end{theorem}
\formalnote{\lean{LanczosRun.gram_commutator_conclusion}}

For a unit vector $y\in\R^k$ satisfying $T_ky=\theta y$, let
$x=Vy$ be its computed Ritz vector.  The last coordinate of $y$ is denoted
$y_k$.

\begin{theorem}[Paige boundary-product bound]\label{thm:paige}
For every unit eigenpair $(\theta,y)$ of $T_k$,
\begin{equation}\label{eq:paige-bound}
  \abs{(x^Tb)y_k}
  \leq 6000k^3\varepsilon\norm H.
\end{equation}
\end{theorem}
\formalnote{\lean{LanczosRun.paige_bound}}

The theorem is an absolute product bound.  It does not separately bound
$\norm b$, $\abs{y_k}$, or the angle between $x$ and $b$.

\section{Ritz-value containment and stabilization}

Define the three nonnegative scales
\begin{align}
 E_k&=6000k^3\varepsilon\norm H,\label{eq:Ek}\\
 F_k&=700k^2\varepsilon\norm H,\label{eq:Fk}\\
 \Psi_k&=8kE_k=48000k^4\varepsilon\norm H,\label{eq:Psik}
\end{align}
and, for $1\leq t\leq k$,
\begin{equation}\label{eq:radius}
 \mathcal R(t)=(8t^2+5)E_k+3F_k.
\end{equation}
The prefix $T_t$ below is the leading $t\times t$ principal submatrix of the
computed tridiagonal $T_k$.

\begin{theorem}[Containment for every prefix]\label{thm:containment}
For every $1\leq t\leq k$ and every unit eigenpair
$T_ty=\theta y$,
\begin{equation}\label{eq:containment}
 \lambda_{\min}(H)-\mathcal R(t)
 \leq\theta\leq
 \lambda_{\max}(H)+\mathcal R(t).
\end{equation}
In particular, this holds with $t=k$ for every eigenvalue of $T_k$.
\end{theorem}
\formalnote{\lean{LanczosRun.containment_prefix}; \lean{LanczosRun.containment}; \lean{LanczosRun.containment_eigenvalue}}

The proof in Lean does not assume a small Gram defect.  It uses a formally
verified descent argument when $\abs{y^TRy}>3/8$.  In the complementary case,
the following sharper local statement is available.

\begin{proposition}[Small-loss localization]\label{prop:small-loss}
If $T_ky=\theta y$, $y^Ty=1$, and $\abs{y^TRy}\leq3/8$, then
\begin{equation}\label{eq:small-loss}
 d_H(\theta)\leq 3\beta_k\abs{y_k}+3F_k.
\end{equation}
\end{proposition}
\formalnote{\lean{LanczosRun.small_loss_localization}}

\begin{theorem}[Stabilized Ritz localization]\label{thm:stabilized}
For every unit eigenpair $T_ky=\theta y$,
\begin{equation}\label{eq:stabilized-max}
 d_H(\theta)
 \leq (k-1)\Psi_k
   +3\max\{\beta_k\abs{y_k},\Psi_k\}+3F_k.
\end{equation}
Consequently,
\begin{equation}\label{eq:stabilized-add}
 d_H(\theta)
 \leq3\beta_k\abs{y_k}+(k+2)\Psi_k+3F_k.
\end{equation}
If either $\beta_k\abs{y_k}\leq\Psi_k$ or $\beta_k=0$, then
\begin{equation}\label{eq:stabilized-terminal}
 d_H(\theta)\leq(k+2)\Psi_k+3F_k.
\end{equation}
The terminal statement also holds directly for every eigenvalue of $T_k$ when
$\beta_k=0$.
\end{theorem}
\formalnote{\lean{LanczosRun.stabilized_max_form}; \lean{LanczosRun.stabilized}; \lean{LanczosRun.stabilized_of_boundary_le}; \lean{LanczosRun.terminal_localization_eigenvalue}}

Because $k>0$, the displayed coefficient $k+2$ is exactly the mathematical
simplification of the formal natural-number expression $(k-1)+3$.

\section{Normalization and the Greenbaum construction}

Let $d_j=\norm{v_j}$ for $1\leq j\leq k$,
$D_v=\diag(d_1,\ldots,d_k)$, and
\[
 \bar V=VD_v^{-1}=[\bar v_1\ \cdots\ \bar v_k].
\]
The run assumptions imply $d_j>0$, so this normalization is well-defined, and
$\norm{\bar v_j}=1$.  If $\beta_k\neq0$, set
\[
 d_{k+1}=\norm{v_{k+1}},\qquad
 \bar v_{k+1}=v_{k+1}/d_{k+1},\qquad
 \bar\beta_k=\beta_kd_{k+1}/d_k.
\]
If $\beta_k=0$, the formal definitions use
$\bar v_{k+1}=\bar v_1$ and $\bar\beta_k=0$.

\begin{theorem}[Normalized local recurrence]\label{thm:normalized-local}
There is a matrix $\bar E\in\R^{n\times k}$ such that
\begin{equation}\label{eq:normalized-recurrence}
 H\bar V=\bar VT_k+\bar\beta_k\bar v_{k+1}e_k^T+\bar E,
\end{equation}
with
\begin{equation}\label{eq:normalized-E}
 \norm{\bar E}\leq2600k^4\varepsilon\norm H.
\end{equation}
For $1\leq j<k$,
\begin{equation}\label{eq:normalized-adjacent}
 \beta_j\abs{\bar v_j^T\bar v_{j+1}}
 \leq4800k\varepsilon\norm H,
\end{equation}
and at the boundary
\begin{equation}\label{eq:normalized-boundary}
 \bar\beta_k\abs{\bar v_k^T\bar v_{k+1}}
 \leq4800k\varepsilon\norm H.
\end{equation}
\end{theorem}
\formalnote{\lean{LanczosRun.normalized_governing}; \lean{LanczosRun.normalized_local_errors}}

Let $U$ be the strictly upper-triangular part of $\bar V^T\bar V$.  Since the
columns of $\bar V$ are unit vectors,
\begin{equation}\label{eq:normalized-gram}
 \bar V^T\bar V=I+U+U^T.
\end{equation}
Set $a=\bar V^T\bar v_{k+1}$ and
\begin{equation}\label{eq:J}
 J=T_kU-UT_k-\bar\beta_k e_ka^T.
\end{equation}

\begin{theorem}[Normalized commutator]\label{thm:normalized-commutator}
The matrix $J$ is upper triangular,
\begin{equation}\label{eq:J-skew}
 J-J^T=\bar V^T\bar E-\bar E^T\bar V,
\end{equation}
and
\begin{equation}\label{eq:J-bound}
 \norm J\leq50000k^5\varepsilon\norm H.
\end{equation}
\end{theorem}
\formalnote{\lean{LanczosRun.normalized_gram_commutator}}

Define the finite geometric transform
\begin{equation}\label{eq:CK}
 C=\sum_{r=0}^{k-1}(-U)^r,
 \qquad K=I-C,
\end{equation}
and set
\begin{equation}\label{eq:sy}
 s=Ca,qquad y_0=\bar v_{k+1}-\bar VCs.
\end{equation}

\begin{theorem}[Nilpotent isometry identities]\label{thm:nilpotent}
The matrix $K$ is strictly upper triangular and
\begin{equation}\label{eq:nilpotent-bounds}
 K^k=0,\qquad \norm K\leq1,\qquad \norm C\leq2.
\end{equation}
Moreover,
\begin{align}
 (\bar VC)^T(\bar VC)&=I-K^TK,\label{eq:nil-gram}\\
 (\bar VC)^Ty_0&=-K^Ts,\label{eq:nil-cross}\\
 \norm s^2+\norm{y_0}^2&=1.\label{eq:nil-norm}
\end{align}
\end{theorem}
\formalnote{\lean{LanczosRun.nilpotent_isometry}}

The notation $\norm s^2$ in \eqref{eq:nil-norm} means
$(\norm s)^2$; both $s$ and $y_0$ use Euclidean norms.

\section{Physical dilation and exact backward model}

Let
\begin{equation}\label{eq:blockH}
 \mathcal H=\diag(H,\ldots,H)\in\R^{nk\times nk}
\end{equation}
contain $k$ copies of $H$.  Form $X\in\R^{nk\times k}$ by stacking the
$k$ blocks
\begin{equation}\label{eq:Xblocks}
 \bar VC,\ \bar VCK,\ \ldots,\ \bar VCK^{k-1}.
\end{equation}
The boundary vector $q\in\R^{nk}$ has first block $y_0$ and, for
$r=1,\ldots,k-1$, block $r+1$ equal to $\bar VCK^{r-1}s$.

\begin{theorem}[Physical dilation]\label{thm:dilation}
The scalar $\bar\beta_k$ is nonnegative and
\begin{equation}\label{eq:dilation-orthogonality}
 X^TX=I,qquad X^Tq=0,qquad \norm q\leq1.
\end{equation}
The first column of $X$ is
$(\bar v_1^T,0,\ldots,0)^T$.  There is
$\widehat E\in\R^{nk\times k}$ such that
\begin{equation}\label{eq:dilation-recurrence}
 \mathcal HX=XT_k+\bar\beta_kq e_k^T+\widehat E,
\end{equation}
and
\begin{equation}\label{eq:dilation-error}
 \norm{\widehat E}
 \leq10^6k^8\varepsilon\norm H.
\end{equation}
At the block level, each of the $k$ propagated recurrence errors is bounded by
$10^6k^7\varepsilon\norm H$.
\end{theorem}
\formalnote{\lean{LanczosRun.physical_dilation}; \lean{LanczosRun.block_recurrence}}

The formal correction is a symmetric matrix
$\widehat\Delta\in\R^{nk\times nk}$ satisfying
\begin{equation}\label{eq:correction}
 \widehat\Delta^T=\widehat\Delta,qquad
 \widehat\Delta X=-\widehat E,qquad
 \norm{\widehat\Delta}\leq3\norm{\widehat E}
 \leq3{,}000{,}000k^8\varepsilon\norm H.
\end{equation}
Consequently, $\widehat H=\mathcal H+\widehat\Delta$ is symmetric and obeys
the exact recurrence
\begin{equation}\label{eq:corrected-recurrence}
 \widehat HX=XT_k+\bar\beta_kq e_k^T.
\end{equation}

\formalnote{\lean{LanczosRun.greenbaumDelta_transpose}; \lean{LanczosRun.greenbaumDelta_mul_X}; \lean{LanczosRun.greenbaumDelta_bound}; \lean{LanczosRun.greenbaum_corrected_recurrence}}

\begin{theorem}[Existence of the Greenbaum model]\label{thm:model}
There exist an integer $d\geq0$, an orthogonal square matrix
\[
 Q\in\R^{nk\times(k+d)},\qquad nk=k+d,qquad Q^TQ=I,quad QQ^T=I,
\]
and a real symmetric tridiagonal matrix
$\widetilde T\in\R^{(k+d)\times(k+d)}$ with the following properties.

\begin{enumerate}[label=(G\arabic*),leftmargin=3.2em]
\item The first $k$ columns of $Q$ are the columns of $X$.
\item The leading $k\times k$ block of $\widetilde T$ is exactly $T_k$.
\item With
\begin{equation}\label{eq:S}
 S=Q^T\mathcal H Q,
\end{equation}
one has
\begin{equation}\label{eq:model-distance}
 \delta:=\norm{S-\widetilde T}
 \leq3\norm{\widehat E}
 \leq3{,}000{,}000k^8\varepsilon\norm H.
\end{equation}
\item The first $k-1$ prescribed subdiagonal entries of $\widetilde T$ are
$\beta_1,\ldots,\beta_{k-1}$ and are positive.
\end{enumerate}
\end{theorem}
\formalnote{\lean{LanczosRun.exists_greenbaumModel}; \lean{LanczosRun.GreenbaumModel.delta_explicit}; \lean{LanczosRun.GreenbaumModel.prescribed_subdiagonal_positive}}

Here ``tridiagonal'' follows from the conjunction of formal symmetry and the
formal upper-tridiagonal predicate.  The theorem does not assert positivity of
the subdiagonal beyond the leading $k\times k$ block.

\begin{theorem}[Exact Lanczos reproduction]\label{thm:reproduction}
Let exact Lanczos be initialized at $e_1\in\R^{k+d}$ and applied to
$\widetilde T$ for $k$ steps.  In the precise sense of the exact-run interface
in Definition~\ref{def:exact-run}, an exact run exists, and every such run
returns the tridiagonal matrix $T_k$.  In particular, its computed
$\alpha_1,\ldots,\alpha_k$ and internal
$\beta_1,\ldots,\beta_{k-1}$ agree with those of the original formal run.
\end{theorem}
\formalnote{\lean{LanczosRun.GreenbaumModel.exactLanczosRun}; \lean{LanczosRun.GreenbaumModel.exact_lanczos_alpha}; \lean{LanczosRun.GreenbaumModel.exact_lanczos_beta}; \lean{LanczosRun.GreenbaumModel.exact_lanczos_returns_Tmat}}

\begin{definition}[Exact-run interface]\label{def:exact-run}
For a symmetric matrix $A$, starting vector $q$, and number of steps $k$, the
formal exact run contains sequences $v_j,\alpha_j,\beta_j$ satisfying
\[
 v_0=0,\quad \beta_0=0,\quad v_1=q,\quad \norm{v_{j+1}}=1
 \quad(0\leq j<k),
\]
\[
 \alpha_{j+1}=v_{j+1}^TAv_{j+1},
\]
and, for $j+1<k$,
\[
 \beta_{j+1}>0,qquad
 \beta_{j+1}v_{j+2}
 =Av_{j+1}-\alpha_{j+1}v_{j+1}-\beta_jv_j.
\]
This is the exact scope of Theorem~\ref{thm:reproduction}; no particular
software implementation of exact Lanczos is invoked.
\end{definition}
\formalnote{\lean{FinitePrecisionLanczos.ExactLanczos}}

\section{Spectral and starting-measure consequences}

Let $q_H=\bar v_1$ and $q_M=e_1\in\R^{k+d}$.  The matrix $S$ in
\eqref{eq:S} is orthogonally similar to $k$ copies of $H$.

\begin{theorem}[Exact starting moments for the unperturbed dilation]
\label{thm:moments}
For every $m\in\N$,
\begin{equation}\label{eq:moments}
 q_M^TS^mq_M=q_H^TH^mq_H.
\end{equation}
Equivalently, the finite atomic starting spectral masses agree exactly:
\begin{equation}\label{eq:mass-equality}
 \mu_{S,q_M}=\mu_{H,q_H}.
\end{equation}
\end{theorem}
\formalnote{\lean{LanczosRun.GreenbaumModel.starting_moment}; \lean{LanczosRun.GreenbaumModel.starting_spectralMass}}

The exact equality is between $H$ and $S$, not between $H$ and
$\widetilde T$.  The latter comparison is perturbative, as in the next two
theorems.

\begin{theorem}[Localization of every model eigenvalue]\label{thm:eigloc}
For every $\widetilde\lambda\in\spec(\widetilde T)$,
\begin{equation}\label{eq:eigloc}
 d_H(\widetilde\lambda)\leq\delta
 \leq3{,}000{,}000k^8\varepsilon\norm H.
\end{equation}
\end{theorem}
\formalnote{\lean{LanczosRun.GreenbaumModel.eigenvalue_localization}; \lean{LanczosRun.GreenbaumModel.eigenvalue_localization_explicit}}

\begin{theorem}[Starting spectral mass in dual Wasserstein distance]
\label{thm:wasserstein}
The physical and model starting masses satisfy
\begin{equation}\label{eq:wasserstein}
 \Wone(\mu_{H,q_H},\mu_{\widetilde T,q_M})
 \leq\sqrt{nk}\,\delta
 \leq\sqrt{nk}\,
       3{,}000{,}000k^8\varepsilon\norm H.
\end{equation}
\end{theorem}
\formalnote{\lean{LanczosRun.GreenbaumModel.starting_wasserstein}; \lean{LanczosRun.GreenbaumModel.starting_wasserstein_explicit}}

We finally state the proved cluster-mass estimate in full.  Let
$\lambda_i$ and $w_i=(u_i^Tq_H)^2$ be the physical eigenvalues and starting
weights.  Let $\widetilde\lambda_j$ and
$\widetilde w_j=(\widetilde u_j^Tq_M)^2$ be their model counterparts.  For a
predicate $G$ on $\R$, define its $\delta$-neighborhood relative to the
physical eigenvalue list by
\begin{equation}\label{eq:cluster-neighborhood}
 C_G=\{t\in\R:\text{there exists }i\text{ with }
     G(\lambda_i)\text{ and }\abs{t-\lambda_i}\leq\delta\}.
\end{equation}

\begin{theorem}[Starting cluster mass]\label{thm:cluster}
Assume that $G$ and membership in $C_G$ are decidable.  Suppose there is
$g\in\R$ such that every pair of physical eigenvalues lying on opposite sides
of $G$ satisfies
\begin{equation}\label{eq:cluster-gap}
 g\leq\abs{\lambda_i-\lambda_j},
\end{equation}
and assume $2\delta<g$.  Then
\begin{equation}\label{eq:cluster-product}
 (g-2\delta)
 \abs{\sum_{i:G(\lambda_i)}w_i-
       \sum_{j:\widetilde\lambda_j\in C_G}\widetilde w_j}
 \leq\sqrt{nk}\,\delta,
\end{equation}
and hence
\begin{equation}\label{eq:cluster}
 \abs{\sum_{i:G(\lambda_i)}w_i-
       \sum_{j:\widetilde\lambda_j\in C_G}\widetilde w_j}
 \leq\frac{\sqrt{nk}\,\delta}{g-2\delta}.
\end{equation}
\end{theorem}
\formalnote{\lean{LanczosRun.GreenbaumModel.starting_cluster_mass_product}; \lean{LanczosRun.GreenbaumModel.starting_cluster_mass}}

The theorem uses the actual model error $\delta$, not merely its explicit
upper bound.  Substituting the bound in \eqref{eq:model-distance} into a
denominator would require separately checking that the resulting denominator
remains positive; that substitution is not packaged as a Lean theorem.

\section{What is, and is not, established}

The verified implication can be summarized as
\[
 \boxed{\text{formal run interface}}
 \Longrightarrow
 \begin{gathered}
 \text{local bounds and exact matrix recurrence}\\
 \text{Paige commutator, loss, containment, stabilization}\\
 \text{normalized recurrence and physical dilation}\\
 \text{symmetric tridiagonal backward model}\\
 \text{exact reproduction and spectral-measure bounds}.
 \end{gathered}
\]

The following claims are intentionally \emph{not} made by the formalization.

\begin{enumerate}[leftmargin=2.3em]
\item It does not derive Assumption~\ref{ass:run} from source code, hardware
floating-point operations, IEEE~754, or a compiler semantics.
\item It does not identify $\varepsilon$ with a particular unit roundoff; the
five local error bounds in Assumption~\ref{ass:run} are hypotheses.
\item It does not assume or prove $k\leq n$.
\item It does not claim exact equality of the starting spectral masses of
$H$ and $\widetilde T$; exact equality holds for $H$ and the unperturbed
dilation $S$, while \eqref{eq:wasserstein} controls $\widetilde T$.
\item It does not claim that all subdiagonal entries of the completed
$\widetilde T$ are positive.  Positivity is verified only for the first
$k-1$ entries inherited from the computed run.
\item The constants are explicit sufficient bounds proved in Lean.  No claim
of sharpness or optimal dependence on $k$ is made.
\end{enumerate}

\section{Declaration correspondence}

Table~\ref{tab:map} is a compact audit map from the displayed mathematical
results to the principal Lean declarations.  Supporting algebraic,
finite-dimensional spectral, tridiagonal-completion, and finite-atomic-measure
lemmas are imported by these declarations but are not individually restated as
journal theorems here.

\begin{table}[htbp]
\centering
\footnotesize
\begin{tabular}{p{0.31\textwidth}p{0.60\textwidth}}
\toprule
Mathematical item & Lean declaration(s) \\
\midrule
Run assumptions & \lean{FinitePrecisionLanczos.LanczosRun} \\
Local errors & \lean{LanczosRun.local_errors} \\
Assembled recurrence & \lean{LanczosRun.governing_matrix}, \lean{LanczosRun.Fmat_bound} \\
Gram commutator & \lean{LanczosRun.gram_commutator_conclusion} \\
Paige bound & \lean{LanczosRun.paige_bound} \\
Prefix containment & \lean{LanczosRun.containment_prefix}, \lean{LanczosRun.containment} \\
Stabilization & \lean{LanczosRun.stabilized_max_form}, \lean{LanczosRun.stabilized} \\
Normalized errors & \lean{LanczosRun.normalized_local_errors} \\
Normalized commutator & \lean{LanczosRun.normalized_gram_commutator} \\
Nilpotent transform & \lean{LanczosRun.nilpotent_isometry} \\
Physical dilation & \lean{LanczosRun.physical_dilation} \\
Symmetric correction & \lean{LanczosRun.greenbaumDelta_bound}, \lean{LanczosRun.greenbaum_corrected_recurrence} \\
Backward model existence & \lean{LanczosRun.exists_greenbaumModel} \\
Model error & \lean{LanczosRun.GreenbaumModel.delta_explicit} \\
Exact reproduction & \lean{LanczosRun.GreenbaumModel.exact_lanczos_returns_Tmat} \\
Starting moments/mass & \lean{LanczosRun.GreenbaumModel.starting_moment}, \lean{LanczosRun.GreenbaumModel.starting_spectralMass} \\
Eigenvalue localization & \lean{LanczosRun.GreenbaumModel.eigenvalue_localization_explicit} \\
Dual Wasserstein bound & \lean{LanczosRun.GreenbaumModel.starting_wasserstein_explicit} \\
Cluster mass & \lean{LanczosRun.GreenbaumModel.starting_cluster_mass} \\
\bottomrule
\end{tabular}
\caption{Principal correspondence between this statement record and Lean.}
\label{tab:map}
\end{table}

\section{Conclusion}

Subject exactly to Assumption~\ref{ass:run}, the formal development verifies a
complete chain from local floating-point-style identities to explicit Paige
and Greenbaum conclusions.  The central backward object is a symmetric
tridiagonal $\widetilde T$ with computed leading block $T_k$ and operator-norm
distance at most
$3{,}000{,}000k^8\varepsilon\norm H$ from an orthogonal copy of $k$ repeated
blocks of $H$.  Exact Lanczos on that model reproduces $T_k$; every model
eigenvalue lies within the same explicit radius of $\spec(H)$; and the
starting spectral measures obey the dual Wasserstein and cluster-mass bounds
stated above.  The formal scope begins at the run interface, rather than at an
executable floating-point algorithm, and that boundary should be preserved in
any comparison with an informal draft.

\end{document}
